\section{Conclusions}
\label{sec:conclusions}

Las cuatro relaciones de herencia (permutación de estados, espejo, reducción de estados y reducción de vecindad) conectan reglas de dominios distintos y reducen las 256 ECA a 81 clases. La forma canónica del grafo de transición distingue con exactitud las 88 clases por permutación y espejo con un solo espacio de 6 células si se usan el grafo completo y su cociente por traslaciones, pero da una comparación binaria y limitada a reglas del mismo dominio.

De los métodos espectrales, los eigenvalores y los momentos espectrales solo ven los atractores: no separan reglas que difieren en sus transitorios, como la 0 y la 8, y se quedan en 77 clases. Los momentos espectrales con masa de atractor conservan la longitud fija de los momentos espectrales y agregan la información de las cuencas: distinguen 86 de las 88 clases (87 junto con el cociente) y separan las reglas 0 y 8, aunque con una distancia pequeña (0.007 con $n = 12$ y $K = 60$) porque la diferencia está en los primeros pasos de un vector que se satura.

Entre dominios distintos, la masa de atractor da distancia exactamente cero a las relaciones que conservan el grafo de transición (permutación de estados, espejo con vecindad simétrica y reducción de vecindad sin desplazamiento) y, en el cociente, también a las que lo cambian por una traslación. Así, la hipótesis de que la posición de los vecinos no cambia el comportamiento se cumple módulo traslaciones, salvo con vecindades no contiguas en espacios de tamaño par. La reducción de estados, en cambio, no se identifica: las configuraciones con el estado nuevo dominan la normalización, y aunque la extensión conserva parcialmente el orden de las distancias (correlación de 0.76 en el cociente), la regla original casi nunca es la más cercana. Los momentos espectrales simples sí son invariantes ante esta relación, de modo que la información de las cuencas se obtiene a costa de esa invariancia.

El vector también depende del tamaño del espacio: la comparación de una regla consigo misma en tamaños distintos solo es confiable para reglas cuyos atractores no dependen de $n$, por lo que conviene comparar en el mismo tamaño. Con esa precaución, la distancia agrupa las ECA en cuatro familias según el periodo dominante de sus atractores (punto fijo, periodo $n$, periodo 2 y periodos largos), aunque las reglas caóticas quedan repartidas porque sus ciclos largos no caben en la ventana.

Como trabajo futuro quedan: normalizar la masa de atractor de modo que sea invariante ante la reducción de estados, por ejemplo restringiéndola a las configuraciones que sobreviven al primer paso o combinándola con los momentos espectrales simples; dar más peso a los primeros pasos y conservar la magnitud del vector, para no perder los transitorios cortos ni los ciclos que no caben en la ventana; elegir de forma sistemática el tamaño del espacio y la ventana $K$, o combinar varios tamaños; y extender la comparación a dominios más grandes mediante muestreo.
